\chapter{Introduction}\label{ch:introduction}

\section{Motivation}\label{sec:motivation}

\emph{Mechanistic interpretability} seeks to reverse-engineer the internal computations of transformer language models into human-understandable terms~\citep{olah2020zoom, elhage2021mathematical}. As these models are deployed in high-stakes settings, developers are increasingly expected to demonstrate that a model is safe, to explain its failures, and to justify relying on it. A central abstraction in this effort is the \emph{circuit}: a sparse subgraph of the model's computational graph, formed by a subset of its attention heads and multi-layer perceptrons (MLPs) and the connections between them, that is responsible for a specific, controlled behavior~\citep{wang2023interpretability}. Identifying the circuit that implements a given behavior is the task of \emph{circuit discovery}. What exactly counts as a circuit, and how its boundaries are drawn, remains contested, and Chapter~\ref{ch:background} treats this question in detail.

Early work identified circuits manually, through painstaking activation patching, until \citet{acdc} introduced an algorithm that automated the search. The field now offers a fast-growing set of such algorithms, each resting on its own assumptions, ablations, and evaluation protocol. This thesis concerns these automated algorithms, and specifically the question of how their outputs can be compared on equal terms.

The state of the art comprises three families, distinguished by how they estimate the importance of each edge: greedy pruning, gradient-based attribution, and gradient-based mask optimization. Automatic Circuit DisCovery (ACDC)~\citep{acdc} prunes greedily, removing edges one at a time and retaining those whose ablation shifts the model's output distribution on the task beyond a chosen threshold. Edge Attribution Patching (EAP)~\citep{syed-etal-2024-attribution} uses gradients to approximate that ablation effect, scoring every edge in a single backward pass at the cost of a first-order approximation error. Edge Pruning~\citep{edgePruning} and DiscoGP~\citep{yu-etal-2025-sheaf} also rely on gradients, but frame discovery as an optimization problem, learning a differentiable mask over the model's edges. These optimization-based methods are faster than greedy pruning and recover substantially sparser circuits that remain faithful to the full model~\citep{edgePruning, yu-etal-2025-sheaf}, which makes them the current frontier of the field.

The output of any single algorithm, however, cannot be treated as ground truth. Two findings explain why. First, there is usually no single ground-truth circuit to recover: the same behavior can be reproduced by many distinct circuits that are all equally faithful to the model, and this non-identifiability is the norm rather than the exception~\citep{meloux2025everything, chen2026rome}. Second, the faithfulness scores used to judge a circuit depend on methodological choices, in particular the ablation used to test it, so a score reflects the experimenter's setup as much as the model itself~\citep{miller2024transformer}. These findings bear directly on what the published evaluations of the optimization-based methods establish. Each method has been compared against prior algorithms under controlled conditions, but in each case under the protocol it was designed for. \citet{edgePruning} evaluate Edge Pruning against ACDC and EAP under resample ablation with a KL-divergence criterion, which is the objective Edge Pruning itself optimizes. \citet{yu-etal-2025-sheaf} evaluate DiscoGP against Edge Pruning, ACDC, and EAP under zero ablation with accuracy-based criteria, which DiscoGP's objective directly optimizes. In both studies the proposed method prevails under its own protocol. Whether these orderings are properties of the algorithms or of the protocols under which they were demonstrated remains untested: no study has varied the ablation applied during discovery and evaluation while holding the algorithms, tasks, and metrics fixed.

\section{Problem Statement}\label{sec:problem}

The motivation above gives rise to two coupled gaps. The first is a tooling gap. The optimization-based methods, Edge Pruning and DiscoGP, are distributed as standalone implementations, each carrying its own evaluation procedure. AutoCircuit~\citep{miller2024transformer}, a standard framework for edge-level circuit evaluation, does not support them. Without a shared pipeline, each of the four algorithms (Edge Pruning, DiscoGP, ACDC, and EAP) can be evaluated only under the protocol built into its own implementation, so the contribution of the algorithm cannot be separated from the contribution of the protocol.

The second is a knowledge gap about the reliability of the optimization-based methods and of the comparisons between them. One open question is the robustness of the reported orderings: each ordering has been established under a single evaluation protocol, and whether it persists when the ablation used during discovery and evaluation changes has not been examined. A second open question is stability. Because a behavior admits many equally faithful circuits, and because the discovery procedure is itself a high-variance statistical estimator, repeated runs of the same method can produce substantially different circuits, so it is unclear how much of a discovered circuit is a stable property of the model rather than an artifact of one particular run. This instability has so far been quantified mainly for attribution-based methods such as EAP, and characterizing it for other discovery methods, including the mask-optimization methods studied here, is explicitly named as an open direction~\citep{meloux2025variance}. A third open question is whether a circuit reflects genuine learned structure rather than a spurious artifact. Interpretability methods are known to return plausible-looking explanations even for randomly initialized networks~\citep{meloux2025deadsalmons}, so a method must be shown to recover structure that the trained model actually contains, and not the kind of pattern it would have produced regardless. Comparing each method against a random-weight null model makes this concrete. None of these questions has been examined for Edge Pruning or DiscoGP under controlled conditions, and closing the tooling gap is the precondition for answering them, since each requires evaluating every method under identical conditions.

\section{Research Questions}\label{sec:rqs}

This thesis integrates Edge Pruning and DiscoGP, in its edge-only configuration, into the AutoCircuit library, enabling a rigorous and controlled comparison against the algorithms the library already provides. The comparison is carried out on small, well-studied transformer models, GPT-2~\citep{gpt2} and a four-layer attention-only model, across three established benchmark tasks: Indirect Object Identification (IOI)~\citep{wang2023interpretability}, Greater-Than~\citep{hanna2023doesgpt2computegreaterthan}, and Docstring completion~\citep{stefanhex_janiak_2023}. The central research question is:

\begin{quote}
  \textbf{How do state-of-the-art optimization-based circuit discovery methods compare against established algorithms when integrated into a common evaluation framework and assessed under identical ablations and metrics?}
\end{quote}

\noindent To address this question systematically, the investigation is structured around three sub-questions:

\begin{description}
  \item[Sub-RQ 1 (Integration).] How can Edge Pruning and DiscoGP be integrated into the edge-centric architecture of the AutoCircuit library, so that they are evaluated under the same ablations and metrics as the algorithms it already provides (ACDC and Edge Attribution Patching)?
  \item[Sub-RQ 2 (Faithfulness).] Under identical ablations and metrics, how do the integrated methods compare against the established algorithms (ACDC and EAP) in faithfulness across the benchmark tasks (IOI, Greater-Than, and Docstring), and does the resulting ranking persist when the ablation applied during discovery and evaluation is varied?
  \item[Sub-RQ 3 (Reliability).] How reliable are the resulting circuits, in terms of their stability across repeated runs and whether they reflect genuine learned structure rather than spurious artifacts, tested against a random-weight null model?
\end{description}

\section{Contributions}\label{sec:contributions}

This thesis makes the following contributions:

\begin{itemize}
  \item A unified, open-source integration of Edge Pruning and DiscoGP into the AutoCircuit framework, the first setting in which the ablation applied during the discovery and evaluation of these optimization-based methods can be varied independently of the algorithm.
  \item A controlled comparison of all four algorithms across three benchmark tasks that crosses discovery and evaluation ablations systematically, testing whether the orderings reported in the original evaluations persist outside the protocols under which they were established.
  \item An empirical characterization of the reliability of the discovered circuits, in terms of their stability across repeated runs and their behavior on a random-weight null model, neither of which has previously been examined for the optimization-based methods.
\end{itemize}

\section{Thesis Outline}\label{sec:outline}

The remainder of this thesis is organized as follows. Chapter~\ref{ch:background} introduces the necessary background on transformers, the circuit abstraction, and the AutoCircuit framework. Chapter~\ref{ch:related-work} surveys related work on circuit discovery algorithms and on the reliability of circuit faithfulness. Chapter~\ref{ch:methodology} describes the integration of Edge Pruning and DiscoGP into AutoCircuit and the design of the comparison experiments. Chapter~\ref{ch:experiments} details the benchmark tasks, metrics, and experimental setup, and Chapter~\ref{ch:results} reports the comparison results on faithfulness, stability, and the random-weight null test. Chapter~\ref{ch:discussion} interprets these findings and discusses their limitations, and Chapter~\ref{ch:conclusion} concludes and outlines directions for future work.
